\documentclass[12pt, a4paper]{article}
\usepackage[utf8]{inputenc}
\usepackage[T1]{fontenc}
\usepackage[ngerman]{babel}
\usepackage{amsmath, amssymb, amsthm}
\usepackage{geometry}
\usepackage{graphicx}
\usepackage{xcolor}
\usepackage{tcolorbox}
\usepackage{enumitem}
\usepackage{listings}
\usepackage{fancyhdr}
\usepackage{titlesec}
\usepackage{mdframed}
\usepackage{tikz}
\usepackage{pgfplots}
\pgfplotsset{compat=1.18}

\geometry{a4paper, margin=2.5cm}

\definecolor{theoryblue}{RGB}{30, 90, 160}
\definecolor{exerciseorange}{RGB}{200, 80, 0}
\definecolor{solutiongreen}{RGB}{0, 120, 60}
\definecolor{hintgray}{RGB}{100, 100, 100}
\definecolor{codebg}{RGB}{245, 245, 245}

% Theory box
\tcbuselibrary{skins, breakable}
\newtcolorbox{theorybox}[1]{
  colback=theoryblue!5,
  colframe=theoryblue,
  fonttitle=\bfseries\color{white},
  title=#1,
  breakable
}

% Exercise box
\newtcolorbox{exercisebox}[1]{
  colback=exerciseorange!5,
  colframe=exerciseorange,
  fonttitle=\bfseries\color{white},
  title=Aufgabe #1,
  breakable
}

% Solution space box
\newtcolorbox{solutionspace}[1][6cm]{
  colback=white,
  colframe=black!20,
  fonttitle=\itshape\color{hintgray},
  title=Rechenraum,
  height=#1,
  breakable
}

% Hint box
\newtcolorbox{hintbox}{
  colback=solutiongreen!5,
  colframe=solutiongreen,
  fonttitle=\bfseries\color{white},
  title=Hinweis,
  breakable
}

% Julia code style
\lstdefinelanguage{Julia}{
  keywords={function, end, return, if, else, elseif, for, while, in, begin, let, do, try, catch, finally, using, import, export, mutable, struct, abstract, type},
  keywordstyle=\color{theoryblue}\bfseries,
  comment=[l]{\#},
  commentstyle=\color{hintgray}\itshape,
  stringstyle=\color{solutiongreen},
  morestring=[b]",
  sensitive=true
}
\lstset{
  language=Julia,
  backgroundcolor=\color{codebg},
  basicstyle=\ttfamily\small,
  breaklines=true,
  frame=single,
  rulecolor=\color{black!30},
  numbers=left,
  numberstyle=\tiny\color{hintgray}
}

\pagestyle{fancy}
\fancyhf{}
\rhead{\textcolor{hintgray}{\small Rocket Science Workbook}}
\lhead{\textcolor{hintgray}{\small \leftmark}}
\cfoot{\thepage}

\titleformat{\section}{\large\bfseries\color{theoryblue}}{}{0em}{}[\titlerule]
\titleformat{\subsection}{\normalsize\bfseries\color{exerciseorange}}{}{0em}{}


\begin{document}

\begin{center}
  {\LARGE \textbf{Kapitel 01: Das Euler-Verfahren}}\\[0.5em]
  {\large \textcolor{hintgray}{Numerische Lösung gewöhnlicher Differentialgleichungen}}\\[0.3em]
  {\small \textcolor{hintgray}{Rocket Science Workbook — simulation-2D-jl}}
\end{center}
\vspace{1em}
\hrule
\vspace{1.5em}

% ─────────────────────────────────────────────
\section{Was ist eine Differentialgleichung?}

\begin{theorybox}{Definition: Gewöhnliche DGL (ODE)}
Eine \textbf{gewöhnliche Differentialgleichung} (ODE) beschreibt, wie sich eine Größe $y(t)$ \textit{in Abhängigkeit von sich selbst und der Zeit} verändert:
\[
  \frac{dy}{dt} = f(t,\, y)
\]
$f(t,y)$ ist eine beliebige Funktion. Der Wert $y(0) = y_0$ heißt \textbf{Anfangsbedingung}.
Das Ziel: Finde $y(t)$ für alle $t > 0$.
\end{theorybox}

\subsection*{Warum numerisch?}
Die meisten ODEs in der Physik haben \textbf{keine geschlossene analytische Lösung}.
Stattdessen berechnen wir \textit{Näherungswerte} auf einem Zeitgitter:
\[
  t_0,\; t_1 = t_0 + h,\; t_2 = t_0 + 2h,\; \ldots
\]
wobei $h$ die \textbf{Schrittweite} ist.

% ─────────────────────────────────────────────
\section{Das Euler-Verfahren}

\begin{theorybox}{Formel: Explizites Euler-Verfahren}
\[
  \boxed{y_{n+1} = y_n + h \cdot f(t_n,\, y_n)}
\]
\textbf{Idee:} Wir linearisieren $y(t)$ lokal — die Ableitung $f(t_n, y_n)$ wird als konstant über den ganzen Schritt angenommen.
\end{theorybox}

\subsection*{Geometrische Interpretation}
Am Punkt $(t_n, y_n)$ berechnen wir die Steigung der Kurve. Dann gehen wir einen Schritt $h$ entlang dieser Steigung. Je kleiner $h$, desto genauer.

\begin{theorybox}{Fehleranalyse}
\begin{itemize}
  \item \textbf{Lokaler Fehler} (pro Schritt): $\mathcal{O}(h^2)$ — quadratisch klein
  \item \textbf{Globaler Fehler} (nach $N$ Schritten): $\mathcal{O}(h)$ — linear in $h$
  \item Das Euler-Verfahren ist ein Verfahren \textbf{erster Ordnung}
\end{itemize}
Halbiert man $h$, halbiert sich auch der globale Fehler.
\end{theorybox}

% ─────────────────────────────────────────────
\section{Rechenaufgaben}

\begin{exercisebox}{1 — Euler von Hand (5 Schritte)}
Gegeben sei die ODE:
\[
  \frac{dy}{dt} = -y, \quad y(0) = 1, \quad h = 0{,}2
\]
\textbf{Exakte Lösung:} $y(t) = e^{-t}$

Berechne die Näherungswerte $y_1, y_2, y_3, y_4, y_5$ mit dem Euler-Verfahren.
Fülle die Tabelle aus:

\begin{center}
\begin{tabular}{|c|c|c|c|c|}
\hline
$n$ & $t_n$ & $y_n$ (Euler) & $y(t_n) = e^{-t_n}$ & Fehler $|y_n - y(t_n)|$ \\
\hline
0 & 0{,}0 & 1{,}0000 & 1{,}0000 & 0{,}0000 \\
\hline
1 & 0{,}2 & & & \\
\hline
2 & 0{,}4 & & & \\
\hline
3 & 0{,}6 & & & \\
\hline
4 & 0{,}8 & & & \\
\hline
5 & 1{,}0 & & & \\
\hline
\end{tabular}
\end{center}
\end{exercisebox}

\begin{solutionspace}[7cm]
\end{solutionspace}

\begin{hintbox}
Schritt 1: $y_1 = y_0 + h \cdot f(t_0, y_0) = 1{,}0 + 0{,}2 \cdot (-1{,}0) = 0{,}8$\\
Schritt 2: $y_2 = y_1 + h \cdot f(t_1, y_1) = 0{,}8 + 0{,}2 \cdot (-0{,}8) = ?$\\
\ldots weiter bis $n=5$.
\end{hintbox}

\vspace{1em}

\begin{exercisebox}{2 — Einfluss der Schrittweite}
Wiederhole Aufgabe 1 mit $h = 0{,}5$ (nur 2 Schritte bis $t=1$).\\[0.5em]
\textbf{Frage:} Wie groß ist der Fehler bei $t=1{,}0$ im Vergleich zu $h=0{,}2$?\\
Stimmt die erwartete lineare Fehlerordnung $\mathcal{O}(h)$?

\begin{center}
\begin{tabular}{|c|c|c|c|c|}
\hline
$n$ & $t_n$ & $y_n$ (Euler, $h=0{,}5$) & $y(t_n) = e^{-t_n}$ & Fehler \\
\hline
0 & 0{,}0 & 1{,}0000 & 1{,}0000 & 0{,}0000 \\
\hline
1 & 0{,}5 & & & \\
\hline
2 & 1{,}0 & & & \\
\hline
\end{tabular}
\end{center}

Verhältnis der Fehler: $\dfrac{\text{Fehler}_{h=0{,}5}}{\text{Fehler}_{h=0{,}2}} \approx$ \underline{\hspace{2cm}} \quad (erwartet: $\approx 2{,}5$)
\end{exercisebox}

\begin{solutionspace}[5cm]
\end{solutionspace}

\vspace{1em}

\begin{exercisebox}{3 — Freier Fall als ODE}
Ein Objekt fällt aus $h_0 = 100\,\text{m}$ mit $v_0 = 0\,\text{m/s}$. Keine Luftreibung.\\[0.5em]
Das System lautet (Zustandsvektor $\mathbf{y} = [h, v]^\top$):
\[
  \frac{d}{dt}\begin{pmatrix} h \\ v \end{pmatrix}
  = \begin{pmatrix} v \\ -g \end{pmatrix}, \quad g = 9{,}81\,\frac{\text{m}}{\text{s}^2}
\]

Berechne mit $h_{\text{step}} = 1\,\text{s}$ die ersten \textbf{3 Schritte}:

\begin{center}
\begin{tabular}{|c|c|c|c|c|}
\hline
$n$ & $t$ (s) & $h$ (m) & $v$ (m/s) & $v_{\text{exakt}} = -g\cdot t$ \\
\hline
0 & 0 & 100{,}00 & 0{,}00 & 0{,}00 \\
\hline
1 & 1 & & & \\
\hline
2 & 2 & & & \\
\hline
3 & 3 & & & \\
\hline
\end{tabular}
\end{center}

\textbf{Wann trifft das Objekt den Boden?} Exakt: $t^* = \sqrt{2h_0/g} \approx$ \underline{\hspace{2cm}} s
\end{exercisebox}

\begin{solutionspace}[6cm]
\end{solutionspace}

\begin{exercisebox}{4 — Konzeptfrage: Stabilität}
Betrachte die ODE $\dfrac{dy}{dt} = \lambda y$ mit $\lambda > 0$ (wachsende Lösung).\\[0.5em]
\textbf{a)} Schreibe die Euler-Formel für diesen Fall auf.\\[0.3em]
\textbf{b)} Ab welcher Schrittweite $h$ wird $|y_{n+1}| > |y_n|$ garantiert?\\[0.3em]
\textbf{c)} Ab welchem $h$ ``explodiert'' die Näherung (Instabilität)?\\
\textit{Hinweis: Betrachte den Faktor $(1 + h\lambda)$ und wann dieser $> 1$ oder $< -1$ wird.}
\end{exercisebox}

\begin{solutionspace}[6cm]
\end{solutionspace}

% ─────────────────────────────────────────────
\section{Verbindung zum Julia-Template}

\begin{theorybox}{Was du jetzt implementierst}
In \texttt{template.jl} (Kapitel 01) implementierst du:
\begin{enumerate}
  \item \texttt{euler\_step(f, t, y, h)} — einen einzelnen Euler-Schritt
  \item \texttt{euler\_solve(f, y0, t0, t\_end, h)} — die vollständige Integration
  \item Test mit $\dot{y} = -y$ und Fehlerplot
  \item Test mit freiem Fall (Aufgabe 3)
\end{enumerate}
Die Funktion \texttt{euler\_step} ist exakt die Formel $y_{n+1} = y_n + h \cdot f(t_n, y_n)$, die du oben von Hand angewendet hast.
\end{theorybox}

\end{document}
